\documentclass[10pt,a4paper]{amsart}
\usepackage[margin=19mm]{geometry}
\usepackage{amsmath,amssymb,mathtools}
\usepackage[scr=rsfs,cal=euler]{mathalfa}
\usepackage{xfrac}
\usepackage[unicode,hidelinks,bookmarks=false]{hyperref}
\newcommand{\ie}{i.e.\ }
\newcommand{\cf}{cf.\ }
\newcommand{\resp}{respectively\ }
\newcommand{\Subsection}{\subsection}
\providecommand{\nobreakdash}{\nobreak\hbox{-}}
\setlength{\emergencystretch}{2em}
\multlinegap0pt
\usepackage[shortlabels]{enumitem}
\usepackage{amssymb}
\usepackage{mathtools}
\newtheorem{lemma}{Lemma}
\newcommand{\Hess}{\operatorname{Hess}}
\newcommand{\dist}{\operatorname{dist}}
\title{Almost Positive Sectional Curvature with Bounded Geometry on $S^2\times S^2$}
\author{Yuhang Liu}
\date{Source: arXiv:2608.20888v1 (CC BY 4.0)}
\begin{document}
\maketitle

\noindent\emph{Source and licence.} Almost Positive Sectional Curvature with Bounded Geometry on $S^2\times S^2$. Version arXiv:2608.20888v1 (CC BY 4.0), distributed by arXiv under the Creative Commons Attribution 4.0 International licence, http://creativecommons.org/licenses/by/4.0/, as recorded in the arXiv licence metadata for this exact version and shown on the abstract page https://arxiv.org/abs/2608.20888v1. Full licence notice: \texttt{licenses/arxiv-2608.20888-CC-BY-4.0.txt}.\par\smallskip

\noindent\textit{Editorial scope.} Counted unit: the article's lemma lem:needle (source0.tex lines 263--289). The article's complete original proof of it is delivered with it (source0.tex lines 291--379, one \texttt{proof} environment of 89 lines). No mathematics is written, repaired or re-derived: the statement and the proof are the article's own lines, in the article's own notation, and no retained line is substituted or re-flowed.  No further source-local context is needed: every cross-reference of every retained line resolves inside the two delivered spans.  Nothing was omitted from any retained span, so the package carries no omission record.  No retained line contains a citation, so the wrapper carries no bibliography and no bibliography entry.  No retained line contains a figure environment, an image-inclusion command or drawing code, so no image is delivered and the package declares no asset dependency.  Editorial formatting only, in addition: an AMS \texttt{amsart} preamble that restates the article's own declarations and macros used by the retained lines, namely the environments \texttt{lemma} on separate counters, together with the standard packages those lines need.  The retained environments print in this document as Lemma 1 in order of appearance, whereas the article prints the same items with the numbering of its own full document.  The complete native source of the article as distributed in the arXiv e-print for this exact version is bundled unchanged as \texttt{sources/208211/source0.tex}.
\begin{lemma}[local needle]\label{lem:needle}
Let $(N,g)$ be a smooth closed Riemannian manifold and let $q\in N$.  There are constants
\[
 R_*>0,\qquad c_*>0,\qquad C_*>1,
\]
depending only on $(N,g,q)$, with the following property.  For every
$0<\sigma\le1$ and every $0<R\le R_*$ there is a smooth function
$w=w_{\sigma,R}:N\to[0,\infty)$ such that
\begin{enumerate}[label=\textup{(\roman*)}]
 \item $w=0$ outside $B_g(q,R)$ and in a neighborhood of its boundary;
 \item
 \[
  -C_*g\le\Hess_g w\le\sigma g;
 \]
 \item for some $\rho\ge c_*\sigma R$,
 \[
  \Hess_g w\le-3g
  \quad\text{on }B_g(q,\rho);
 \]
 \item
 \[
  |dw|_g\le C_*\sigma R,
  \qquad
  0\le w\le C_*\sigma R^2.
 \]
\end{enumerate}
\end{lemma}

\begin{proof}
Let $r=\dist_g(q,\cdot)$.  After decreasing $R_*$, geodesic polar coordinates are smooth on $0<r<R_*$ and
\begin{equation}\label{eq:hess-r}
 \Hess_g r
 =\frac1r(g-dr^2)+E_r,
 \qquad
 |E_r|_g\le C r,
\end{equation}
for a fixed constant $C$.  Equivalently,
$\Hess_g(r^2/2)=g+O(r^2)$.

We construct a one-variable function $F$.  Put $A=4$.  Choose smooth functions
$\alpha:[0,\infty)\to[0,1]$ and
$\gamma\in C_c^\infty((0,1))$ such that
\[
 \alpha=1\ \text{on }[0,1],\qquad
 \alpha=0\ \text{on }[2,\infty),\qquad
 \gamma\ge0,\qquad \int_0^1\gamma=1.
\]
Choose a sufficiently small fixed constant $c>0$ and put
$\ell=c\sigma R$.  Define
\[
 n(r)=-A\alpha(r/\ell),
 \qquad
 D=-\int_0^R n(s)\,ds,
\]
and put a positive bump of total mass $D$ in the interval
$(3\ell,R-\ell)$:
\[
 p(r)=
 \frac{D}{R-4\ell}
 \gamma\!\left(\frac{r-3\ell}{R-4\ell}\right)
\]
there, and $p=0$ elsewhere.  Taking $c$ small, independently of
$\sigma$ and $R$, ensures
\[
 0\le p\le\frac{\sigma}{4},
 \qquad 4\ell<R.
\]
Set
\[
 q_1=n+p,
 \qquad
 Q(r)=\int_0^r q_1(s)\,ds.
\]
The negative bump occurs before the positive one, the two total masses cancel, and therefore
\begin{equation}\label{eq:Q-properties}
 Q\le0,\qquad Q=0\text{ near }R,
 \qquad |Q|\le C_1\sigma R.
\end{equation}
Moreover,
\begin{equation}\label{eq:q-properties}
 -A\le q_1\le\frac{\sigma}{4},
 \qquad
 \frac{Q(r)}r\ge-C_2A.
\end{equation}
Define
\[
 F(r)=-\int_r^R Q(s)\,ds.
\]
Then $F\ge0$, $F'=Q$, $F''=q_1$, and $F=0$ near $R$.  On
$0\le r\le\ell$ one has
\[
 F(r)=F(0)-\frac A2r^2.
\]
The radial function $w=F(r)$, extended by zero, is smooth at both
$q$ and $\partial B_g(q,R)$.

For $r>0$,
\begin{equation}\label{eq:hess-radial-w}
 \Hess_g w=q_1(r)dr^2+Q(r)\Hess_g r.
\end{equation}
The main tangential term $Q(r)r^{-1}(g-dr^2)$ is nonpositive.  The error in \eqref{eq:hess-r} has norm at most
$C|Q|r\le C C_1\sigma R^2$.  Decrease $R_*$ so that this is at most $3\sigma/4$.  Equations
\eqref{eq:q-properties} and \eqref{eq:hess-radial-w} then give
$\Hess w\le\sigma g$, while their lower bounds give
$\Hess w\ge-C_*g$ for a fixed $C_*$.  On $r\le\ell$,
\[
 \Hess_gw=-A\Hess_g(r^2/2).
\]
After one more decrease of $R_*$, this is at most $-3g$.
Thus one may take $\rho=\ell=c\sigma R$.  Finally,
\eqref{eq:Q-properties} gives
\[
 |dw|=|Q|\le C_*\sigma R,
 \qquad
 0\le w\le R\max|Q|\le C_*\sigma R^2.
\]
\end{proof}

\end{document}
